% GENERATED FILE. Edit canonical lessons, then run npm run generate:books.
% canonical-source-hash: fnv1a64:1f3358e3a04c68cd
% canonical-lessons: PA-R164-wake-again, PA-R164-drive-again, PA-R164-say-again

\chapter{Return to Waking, Driving, and Saying}
\label{ch:pa-three-familiar-actions-r4}
\hlchaptermodality{164}

\begin{chapteropening}
\textbf{By the end of this chapter:} \emph{I can retrieve three already-taught Punjabi action verbs after a long gap, one at a time, without claiming scored A1 writing or complete mock credit.}

Recall waking, driving, and saying from short familiar cues without introducing new Punjabi.
\end{chapteropening}

\section[\texorpdfstring{\pa{ਜਾਗਣਾ} (jāgṇā)}{ਜਾਗਣਾ (jāgṇā)}]{An old word for waking}

\label{lesson:PA-R164-wake-again}



\begin{quote}
Hear \textbf{to wake up}. Recall \textbf{\pa{ਜਾਗਣਾ}} (\emph{jāgṇā}).
\end{quote}

\subsection*{Your turn}
Hear \textbf{\pa{ਜਾਗਣਾ}}. Recall \textbf{to wake up}. Now hear \textbf{to wake up} again. Say \textbf{\pa{ਜਾਗਣਾ}} before the model. One old word is enough.

\begin{tcolorbox}[breakable,colback=teal!4,colframe=teal!35!black,title={Before you move on}]
One last cue: \textbf{to wake up}. Recall \textbf{\pa{ਜਾਗਣਾ}}. Listen and retry once if the word does not come back.
\end{tcolorbox}
\section[\texorpdfstring{\pa{ਚਲਾਉਣਾ} (calāuṇā)}{ਚਲਾਉਣਾ (calāuṇā)}]{An old word for driving}

\label{lesson:PA-R164-drive-again}



\begin{quote}
Hear \textbf{to drive}. Say \textbf{\pa{ਚਲਾਉਣਾ}} (\emph{calāuṇā}) before the model.
\end{quote}

\subsection*{Your turn}
Hear \textbf{\pa{ਚਲਾਉਣਾ}}. Recall \textbf{to drive}. Then hear \textbf{to drive} again. Recall \textbf{\pa{ਚਲਾਉਣਾ}}. Only this old meaning and word are needed.

\begin{tcolorbox}[breakable,colback=teal!4,colframe=teal!35!black,title={Before you move on}]
Hear \textbf{to drive} once more. Say \textbf{\pa{ਚਲਾਉਣਾ}}. If needed, listen to the model and try again once.
\end{tcolorbox}
\section[\texorpdfstring{\pa{ਕਹਿਣਾ} (kahiṇā)}{ਕਹਿਣਾ (kahiṇā)}]{An old word for saying}

\label{lesson:PA-R164-say-again}



\begin{quote}
Hear \textbf{to say}. Recall \textbf{\pa{ਕਹਿਣਾ}} (\emph{kahiṇā}).
\end{quote}

\subsection*{Your turn}
Pause after each familiar meaning, then hear its model:

\begin{itemize}
\raggedright
  \item \textbf{to drive} — \textbf{\pa{ਚਲਾਉਣਾ}} (\emph{calāuṇā}).
  \item \textbf{to say} — \textbf{\pa{ਕਹਿਣਾ}} (\emph{kahiṇā}).
  \item \textbf{to wake up} — \textbf{\pa{ਜਾਗਣਾ}} (\emph{jāgṇā}).
\end{itemize}

Only the order of these old words has changed.

\begin{tcolorbox}[breakable,colback=teal!4,colframe=teal!35!black,title={Before you move on}]
Hear \textbf{to wake up}, \textbf{to say}, then \textbf{to drive}. Pause after each cue and recall its Punjabi word before the model: \textbf{\pa{ਜਾਗਣਾ}}, \textbf{\pa{ਕਹਿਣਾ}}, \textbf{\pa{ਚਲਾਉਣਾ}}. Revisit only the old word you missed.
\end{tcolorbox}
